\documentclass[12pt,a4paper]{article}
\usepackage[margin=1in]{geometry}
\usepackage{enumitem}
\usepackage{booktabs}
\usepackage{longtable}
\usepackage{hyperref}
\usepackage{xcolor}

\hypersetup{colorlinks=true, linkcolor=blue!60!black, urlcolor=blue!60!black}

\title{Paper Status Assessment and Recommendations\\[6pt]
\large ``Quantity-Space Prior Specification for Bayesian Hierarchical Models''}
\author{Technical Note}
\date{February 28, 2026}

\begin{document}
\maketitle

\section{Overview}

This note is a follow-up to the assessment of February 25, 2026. It provides
an updated status of all issues identified in that note, flags new issues
discovered during this review, and refines the recommended order of operations.

The assessment is based on a full read of \texttt{doc/main.tex} (1748 lines),
\texttt{TODO.md}, and the February 25 technical note as of February 28, 2026.

\subsection{Summary of Changes Since February 25}

No substantive changes have been made to \texttt{doc/main.tex} since the
February 25 assessment. All Priority~1 items (immediate text fixes) were
completed during that session. The Priority~2 structural improvements,
the simulation scope decision (Priority~3), and the posterior sensitivity
analysis (Priority~4) remain open.

On the code side, Phase~3 (CE-space priors) remains in progress: the CE
Jacobian, CE-space posterior, CE-space prior selection, and CE-space sampling
modules are implemented with Normal population distributions, but the
Beta-Beta hierarchy for bounded CEs is still TODO. Phase~2 ($\kappa$
parameterization and sensitivity grid) is complete.


\section{Issues Carried Forward from February 25}

The following table updates the status of all issues from the previous note.
Items marked ``Fixed'' were resolved during the February 25 session; items
marked ``Open'' remain unaddressed.

\begin{center}
\begin{tabular}{lllll}
\toprule
Location & Type & Severity & Status & Description \\
\midrule
Abstract & Clarity & Minor & Open & ``reduces \ldots to two'' ambiguity ($4p \to 2$ vs $4p \to 2p+2$) \\
Abstract & Scope & Minor & Open & ``any distributional family'' claim \\
\S2.2 (303--338) & Structure & Moderate & Open & Belief-elicitation detour \\
\S3 title & Structure & Minor & Open & Title doesn't reflect both contributions \\
\S3 ``no free lunch'' & Placement & Minor & Open & Paragraph at end of \S3.3 feels like afterthought \\
\S4.1 (660--662) & Clarity & Minor & Open & Wilcox (2011) normalization unexplained \\
\S4.2 (668--675) & Style & Minor & \textbf{New} & ``we'' $\to$ ``I'' (3 instances missed in Feb 25) \\
\S4.3 & Length & Minor & Open & Injectivity could be condensed \\
\S4.7 & Missing & Moderate & Open & No posterior sensitivity analysis \\
\S7.3 & Argument & Minor & Open & BF ``perverse incentive'' needs citation \\
App.\ E (1663, 1669) & Maintenance & Minor & Open & Hardcoded ``115 blocks'' \\
\S5--6 & Completeness & Major & Open & Stub sections \\
Figures & Captions & Minor & Open & Terse captions (see \S\ref{sec:captions}) \\
Road map (201--213) & Consistency & Minor & Open & References stubs without caveat \\
\bottomrule
\end{tabular}
\end{center}


\section{New Issues Found}

\subsection{First-person inconsistencies in \S4.2 (lines 668--675)}

The February 25 note fixed ``we'' on lines 421 and 450 but missed three
instances in \S4.2 (``Quantity-Space Mapping for RDU''):

\begin{enumerate}[label=(\alph*)]
  \item Line 668: ``While we can follow \textbackslash textcite\{Gao2022\}''
    $\to$ ``While one can follow'' or ``While it is possible to follow''
  \item Line 673: ``we start instead by specifying'' $\to$ ``I start instead
    by specifying''
  \item Line 675: ``We discuss first'' $\to$ ``I discuss first''
\end{enumerate}

The remaining ``we'' instances (lines 496, 1449, 1451) are in mathematical
exposition (``if we first define,'' ``suppose we wish to specify'') where the
inclusive ``we'' is conventional and acceptable.


\subsection{Figure captions are genuinely terse}
\label{sec:captions}

The \verb|\mycaption| command (defined in \texttt{packages.tex}, line 196) is a
thin wrapper around \verb|\caption| that adjusts paragraph indentation. It adds
no additional content. The actual caption strings in the paper are:

\begin{enumerate}[label=(\alph*)]
  \item ``Hyperparameter posteriors'' (Fig.\ \ref*{} line 960)
  \item ``Hyper-posterior distribution'' (Fig.\ line 966)
  \item ``Hyper-posterior predictive CE distributions'' (Fig.\ line 1014)
  \item ``Posterior predictive CE: Least Affected (Subject \ldots)'' (Fig.\ line 1035)
  \item ``Posterior predictive CE: Most Affected (Subject \ldots)'' (Fig.\ line 1041)
\end{enumerate}

JoE and most economics journals expect self-contained captions that summarize
the key takeaway of the figure. For example, Figure 1 should read something
like: ``Posterior densities of the eight hyperparameters under the informative
prior (blue), diffuse prior (orange), and the informative hyperprior (grey
dashed). The informative hyperprior concentrates posterior mass for all
hyperparameters except $\mu_\lambda$.''

This is a presentation issue, not a methodological one, but it will likely
draw reviewer comments.


\subsection{Road map references subsection label as section}

The road map paragraph (lines 201--213) writes:

\begin{quote}
  ``Section~\verb|\ref{sec:variance_decomposition}| introduces the
  distribution-free variance decomposition\ldots''
\end{quote}

The label \verb|sec:variance_decomposition| is attached to
\verb|\subsection{Hyperprior Construction via Variance Decomposition}| (line
455--456), which is \S3.2. LaTeX will render this as ``Section 3.2,'' which is
technically correct but reads oddly in a road map that otherwise references
top-level sections.

This is harmless if the current structure is retained, but would become a
problem if Section~3 is split (as recommended in Priority~2c of the Feb 25
note). If Section~3 is split, the label should be promoted to a
\verb|\section| and the road map will read naturally.


\subsection{Abstract scope vs demonstrated generality}

The abstract states that the direct prior and $(\rho, \kappa)$ ``fully
determine the hyperprior for any distributional family.'' The theoretical claim
is correct (moment matching applies to any family with finite first two
moments), but the paper demonstrates this only for the Normal-Inverse-Gamma
family, with the Skew-Normal extension in Appendix~B as a generalization
example. A reviewer may flag the gap between claim and demonstration.

Two options: (a) soften the abstract to ``for standard conjugate families,
with extensions demonstrated for the Skew-Normal''; or (b) leave the claim
and ensure the Skew-Normal appendix is explicitly cited in the abstract as
evidence of generality.


\section{Prioritized Recommendations (Updated)}

\subsection{Priority 1: Text fixes (no computation) --- PARTIALLY COMPLETE}

\begin{enumerate}
  \item[{[x]}] All items from Feb 25 Priority 1 --- complete
  \item[{[ ]}] Fix ``we'' $\to$ ``I'' on lines 668, 673, 675
  \item[{[ ]}] Expand figure captions to be self-contained (5 figures)
  \item[{[ ]}] Add one-sentence clarification of Wilcox (2011) normalization
    in \S4.1 (or forward reference to Appendix E, line 1711)
  \item[{[ ]}] Make ``115 blocks'' in Appendix E dynamic via
    \verb|\py{int(Q.get('n_subjects', 0)) + 4}|
\end{enumerate}

These are all quick fixes that can be done in a single editing session.


\subsection{Priority 2: Structural improvements --- UNCHANGED}

These remain as recommended in the February 25 note:

\begin{enumerate}[label=(\alph*)]
  \item Condense belief-elicitation literature in \S2.2 to 2--3 sentences or a
    footnote. Saves approximately 1 page.
  \item Condense main-text injectivity discussion (\S4.3). Keep the
    proposition and the practical recommendation; move compensation mechanism
    details to Appendix~C.
  \item Consider splitting Section~3 into two sections. If done, promote
    \verb|sec:variance_decomposition| to a \verb|\section| label. The road map
    paragraph would then reference full sections throughout.
  \item Strengthen or shorten the Bayes factor discussion (\S7.3). Either add
    a citation (Gelman \& Shalizi 2013, or the ``garden of forking paths''
    literature) or shorten to a brief acknowledgment.
\end{enumerate}


\subsection{Priority 3: Decide simulation study scope --- DECISION STILL NEEDED}

This remains the most consequential decision. The three options from the
February 25 note are unchanged:

\begin{enumerate}[label=(\alph*)]
  \item \textbf{Option A (recommended): CPT simulation only.} Drop discounting
    or reduce to a brief discussion. CPT has the strongest $4p$ motivation
    (8--10 parameters, 32--40 hyperprior choices). Discounting can be added in
    revision.
  \item \textbf{Option B: Both simulations.} Retain current structure. Requires
    implementing two model posteriors, DGP code, and hundreds of MCMC
    replications.
  \item \textbf{Option C: Real data only.} Drop both stubs. Reframe as a
    methods paper with a single worked application. Requires strengthening the
    sensitivity analysis to compensate.
\end{enumerate}

The recommendation remains Option~A. If this decision is made, the stub
Section~7 (Discounting) and associated Appendix~F content can be removed, and
the paper drops from 9 to 8 sections. The road map, introduction (line 196:
``Two simulation studies''), and conclusion would need corresponding updates.


\subsection{Priority 4: Posterior sensitivity analysis --- UNCHANGED}

The current sensitivity table (Table~1) shows prior predictive invariance
across the $(\rho, \kappa)$ grid, which is by construction. A posterior
sensitivity table would demonstrate practical relevance.

Two approaches remain viable:

\begin{enumerate}[label=(\alph*)]
  \item \textbf{Full MCMC grid:} Re-estimate at each $(\rho, \kappa)$ point.
    Definitive but expensive. A reduced grid (e.g., 4 points: corners of the
    $\rho \times \kappa$ space) would be a practical compromise.
  \item \textbf{Analytical approximation:} For a Normal-Normal model, the
    posterior mean is a precision-weighted average of prior and data. This gives
    closed-form expressions for how posterior summaries shift with $(\rho,
    \kappa)$, without re-running MCMC. The approximation quality depends on how
    well the Normal-Normal model captures the actual posterior geometry.
\end{enumerate}

A third approach worth considering: \textbf{importance sampling reweighting.}
The existing MCMC samples from the baseline $(\rho, \kappa)$ can be
reweighted to approximate the posterior under alternative $(\rho, \kappa)$
values, since only the hyperprior changes. This avoids new MCMC runs entirely
and is exact up to the quality of the importance weights. It works well when
the alternative hyperprior is not too different from the baseline, which is
likely the case for moderate changes in $(\rho, \kappa)$.


\subsection{Priority 5: CE-space implementation (Phase 3)}

Phase 3 status per \texttt{TODO.md}:

\begin{enumerate}
  \item[{[x]}] CE Jacobian via \texttt{jax.jacfwd}
  \item[{[x]}] CE-space posterior (\texttt{rdu\_bhm\_ce.py})
  \item[{[x]}] CE-space prior selection
  \item[{[x]}] CE-space sampling module
  \item[{[x]}] Smoke test (10 iter, 40 chains)
  \item[{[ ]}] Beta-Beta hierarchical moment-matching (code)
  \item[{[ ]}] Beta population distribution in \texttt{rdu\_bhm\_ce.py}
  \item[{[ ]}] Full MCMC run with CE-space priors
  \item[{[ ]}] Comparison figures: parameter-space vs CE-space vs diffuse
\end{enumerate}

Completing this phase would allow the paper to demonstrate Contribution~1
empirically. Currently only Contribution~2 (variance decomposition) is
demonstrated with real data. This is important for the paper's credibility:
the CE-space specification is the headline contribution, but the empirical
application uses parameter-space priors.


\subsection{Priority 6: Choice probability prior for $\lambda$ (Phase 4)}

Deferrable to revision. Currently $\lambda$ uses a direct parameter-space
prior, which is adequate for the paper's argument. The block-triangular
Jacobian structure (\S4.2) is already described in the text.


\section{Recommended Order of Operations}

\begin{enumerate}
  \item \textbf{Text fixes} (Priority 1 remaining items): $<$ 1 hour
  \item \textbf{Simulation scope decision} (Priority 3): decision point,
    no computation
  \item \textbf{Structural editing} (Priority 2): half a day
  \item \textbf{Posterior sensitivity} (Priority 4): 1--3 days depending on
    approach (importance reweighting is fastest)
  \item \textbf{CE-space implementation} (Priority 5): multi-day; can proceed
    in parallel with text work
  \item \textbf{Choice probability prior} (Priority 6): defer
  \item \textbf{Simulation study execution} (Priority 3, post-decision): weeks
  \item \textbf{Final polish} (Phase 6): after all above
\end{enumerate}

The simulation scope decision (item 2) gates everything downstream. If
Option~A (CPT only) is selected, the critical path to a submittable draft is:
text fixes $\to$ structural editing $\to$ posterior sensitivity $\to$ CPT
simulation $\to$ CE-space implementation (can overlap with simulation) $\to$
final polish.

If Option~C (real data only) is selected, the path shortens to: text fixes
$\to$ structural editing $\to$ posterior sensitivity $\to$ CE-space
implementation $\to$ final polish. This would produce a submittable draft
substantially faster, at the cost of controlled evidence.


\section{Summary of All Open Issues}

\begin{center}
\small
\begin{tabular}{lllll}
\toprule
Location & Type & Severity & Priority & Description \\
\midrule
\S4.2 (668--675) & Style & Minor & 1 & ``we'' $\to$ ``I'' (3 instances) \\
Figs 1--5 & Captions & Minor & 1 & Terse captions need expansion \\
\S4.1 (660--662) & Clarity & Minor & 1 & Wilcox normalization unexplained \\
App.\ E (1663, 1669) & Maintenance & Minor & 1 & Hardcoded ``115 blocks'' \\
Abstract & Clarity & Minor & 1--2 & ``reduces \ldots to two'' ambiguity \\
Abstract & Scope & Minor & 1--2 & ``any distributional family'' claim \\
\S2.2 (303--338) & Structure & Moderate & 2 & Belief-elicitation detour \\
\S3 title / structure & Structure & Minor & 2 & Section 3 split consideration \\
\S4.3 & Length & Minor & 2 & Injectivity discussion condensable \\
\S7.3 & Argument & Minor & 2 & BF discussion needs citation or trim \\
\S5--6 & Completeness & Major & 3 & Stub sections \\
\S4.7 & Missing & Moderate & 4 & No posterior sensitivity analysis \\
\S4.4--4.6 & Missing & Moderate & 5 & CE-space priors not demonstrated \\
Road map (201--213) & Consistency & Minor & 2--3 & References stubs; subsection as section \\
\bottomrule
\end{tabular}
\end{center}


\end{document}
